\section{Round-five positive closure: exact history, Doob memory, and the area-corrected tangent}
\label{sec:r5-a4}

This section separates the measurable Ray realization from the stronger
model-specific Feller theorem, derives memory decay from the A2 spectral
packet, constructs the Doob eigenfunction on history space, and retains the
antisymmetric deterministic rough-area drift.

\subsection{The canonical history state}

Let \(E\) be the graph-completed collision or physical state and let
\[
 \mathsf H=D(( -\infty,0],E)
\]
with its Borel sigma-field.  Under a stationary prepared path law \(\nu\), set
\(H_t(\omega)=\omega(t+\cdot)|_{(-\infty,0]}\).  A regular conditional law of
\(H_{t+s}\) given \(H_t\) defines a Borel transition kernel
\(K_s^\nu(h,dh')\) for \(\nu\)-almost every history.  Complete the kernel on a
fixed null set by a cemetery history.

Choose a countable algebra \(\mathscr C_0\) of bounded history cylinders which
separates histories modulo equality of all future conditional laws.  For
\(a\in\mathbb Q_{>0}\) put
\[
 R_aF(h)=\int_0^\infty e^{-as}K_s^\nu F(h)\,ds.
\]
Let \(\mathscr R_\nu\) be the rational vector lattice generated by
\(\mathscr C_0\), its resolvents, constants, minima and maxima.  The evaluation
map
\[
 \iota_\nu(h)=(F(h))_{F\in\mathscr R_\nu}
\]
takes values in a countable product of compact intervals.  Its closure after
identifying equal evaluation vectors is denoted \(\mathsf H_\nu^{\rm Ray}\).

\begin{theorem}[Measured Ray realization]
\label{thm:r5-a4-ray}
The space \(\mathsf H_\nu^{\rm Ray}\) is Polish.  Modulo the single completed
\(\nu\)-null exceptional set, its Borel sigma-field agrees with the completed
history sigma-field.  The kernels induce a positive contraction semigroup
\(P_t^\nu\) on \(C_b(\mathsf H_\nu^{\rm Ray})\), strongly continuous for the
bounded-strict topology \(\beta_0\), and the strict graph closure of
\(\mathscr R_\nu\) is a generator core.
\end{theorem}

\begin{proof}
The evaluation closure is a closed subspace of a countable product of compact
metric spaces, hence is Polish.  Separation by \(\mathscr C_0\) and the
monotone-class theorem identify the completed measurable sets.  The resolvent
identity
\[
 R_a-R_b=(b-a)R_aR_b
\]
and positivity show that the transition kernels preserve the Ray lattice and
therefore extend continuously to its evaluation closure.  For
\(F=R_aG\),
\[
 P_tF-F=\int_0^\infty e^{-as}
 (P_{t+s}G-P_sG)\,ds,
\]
and right continuity of the path kernel plus dominated convergence gives
uniform convergence on every compact evaluation set.  Supremum-norm
contraction extends this from the Ray lattice to its \(\beta_0\)-closure.
Finally \(aR_aF\to F\) in \(\beta_0\), and the resolvent equation gives graph
convergence of the corresponding generator approximants.  This proves the
core statement.  No norm-strong assertion on all of \(C_b\) is made.
\end{proof}

\subsection{The explicit exposed-phase history kernel}

For an exposed A2--A3 phase, the fixed cut tower has a normalized
\(g\)-function \(g_\Psi\) with summable variations.  The history kernel is
obtained by adjoining one future symbol according to \(g_\Psi\) and then
flowing through the bounded roof.  On
\[
 d_\beta(h,\tilde h)=
 \sum_{m\ge1}2^{-m}
 (1\wedge d_{J_1}(h|_{[-m,0]},\tilde h|_{[-m,0]}))
\]
let \(\mathrm{BUC}_\beta(\mathsf H)\) be the bounded uniformly continuous
functions.

Let \(\mathcal L_\Psi r_\Psi=e^{\lambda_\Psi}r_\Psi\) and
\(\ell_\Psi\mathcal L_\Psi=e^{\lambda_\Psi}\ell_\Psi\) be the A2 transfer
eigenvectors.  Define the history eigenfunction by the convergent backward
product
\[
 h_\Psi(h)=r_\Psi(h_0)
 \prod_{j\ge1}
 \frac{g_\Psi(h_{-j}\mid h_{-j-1}^{0})}
      {g_0(h_{-j}\mid h_{-j-1}^{0})}e^{-c_\Psi(h_{-j}^{0})},
\]
where \(c_\Psi\) is the normalized one-step pressure increment.  Summable
variations make the logarithm normally convergent.

\begin{theorem}[History Feller and eigenfunction intertwining]
\label{thm:r5-a4-history-feller}
On every compact simple source chart there are constants
\(0<c<C<\infty\) such that
\[
 c\le h_\Psi\le C,
 \qquad P_t^\Psi h_\Psi=e^{\lambda_\Psi t}h_\Psi.
\]
The exposed-phase history semigroup is norm-strong on
\(\mathrm{BUC}_\beta\) and satisfies, for a cylinder depending on the last
\(m\) units,
\[
 \operatorname{Lip}_\beta(P_t^\Psi F)
 \le Ce^{-\gamma(t-m)_+}\operatorname{Lip}_\beta(F)
     +C\|F\|_\infty.
\]
The transfer and history eigenfunctions are related by the coding map, so the
history eigenvalue is exactly the physical pressure root.
\end{theorem}

\begin{proof}
The logarithm of the product defining \(h_\Psi\) has a tail bounded by
\(C\sum_{j\ge n}\rho^j\), uniformly on the source chart.  Positivity on the
finite head and compactness give the two-sided bounds.  Shifting the product
one step telescopes all factors and gives the eigenfunction equation; roof
suspension converts the discrete eigenvalue to \(e^{\lambda_\Psi t}\).

Two histories agreeing for \(n\) symbols have future conditional densities
whose logarithms differ by \(O(\rho^n)\).  Couple them until their next common
magnet return.  The exponential return tail gives the displayed Lipschitz
estimate.  It implies norm continuity on cylinders and, by density and
contraction, on \(\mathrm{BUC}_\beta\).  The first coordinate of the product is
the A2 right eigenfunction and the remaining factors are the stable-history
Jacobian, proving the intertwining assertion.
\end{proof}

\subsection{Domain-safe finite-dimensional memory}

Let \(U_t\) be the prepared Koopman group on \(L^2(\nu)\).  Let
\(V\subset D(L^r)\), \(r\ge4\), be finite dimensional, let \(P\) be the
orthogonal projection onto \(V\), and write
\[
 C(t)=PU_tP|_V,\qquad A=C'(0)=PLP|_V.
\]
The graph norms of the chosen basis are included in every constant below.

\begin{theorem}[Volterra kernel without formal orthogonal dynamics]
\label{thm:r5-a4-volterra}
There is a unique continuous matrix kernel \(K:[0,\infty)\to\operatorname{End}V\)
such that
\[
 C'(t)=AC(t)+\int_0^tK(t-s)C(s)\,ds.
\]
For \(B\in D(L)\), the resolved coordinate
\(x_B(t)=PU_tB\) satisfies
\[
 x_B'(t)=Ax_B(t)+\int_0^tK(t-s)x_B(s)\,ds+\eta_B(t),
\]
where \(\eta_B\) is continuous and orthogonal to \(V\) at time zero.  If
\(V\subset\widetilde V\), the two kernels are related by the exact block
Schur-complement identity; they are not asserted to be equal.
\end{theorem}

\begin{proof}
Since \(V\subset D(L^2)\), the matrix \(C\) is \(C^2\) and \(C(0)=I\).
Put \(G(t)=C''(t)-AC'(t)\).  Differentiating the desired equation gives the
second-kind Volterra equation
\[
 K(t)+\int_0^tK(t-s)C'(s)\,ds=G(t).
\]
Successive approximation converges uniformly on compact intervals and gives a
unique continuous solution.  Integrating back and using the value at zero
proves the first equation.  Apply it to the cross-correlation
\(PU_tB\) and define the residual as the difference; its continuity follows
from \(B\in D(L)\) and local boundedness of the convolution.  Block inversion
of \(P_{\widetilde V}(z-L)^{-1}P_{\widetilde V}\) gives the stated Schur
complement under enlargement.
\end{proof}

\subsection{Model-derived exponential memory decay}

For \(\Re z>0\),
\[
 \widehat C(z)=P(z-L)^{-1}P,
 \qquad
 \widehat K(z)=zI-A-\widehat C(z)^{-1}.
\]
A2's cut-tower resolvent and high-frequency estimate apply to all matrix
entries because \(V\subset D(L^r)\).  On the mean-zero resolved space, the
cohomological covariance theorem excludes a zero mode.

\begin{theorem}[Meromorphic continuation and decaying memory]
\label{thm:r5-a4-memory-decay}
There is \(\gamma>0\) such that \(\widehat C\) extends meromorphically to
\(\Re z> -\gamma\), has no pole or zero in this strip on the mean-zero
resolved space, and satisfies
\[
 \|\partial_z^j\widehat C(z)\|
 \le C_j(1+|\Im z|)^{-r+j+A},
 \qquad 0\le j\le r-2.
\]
Consequently \(\widehat K\) is analytic in a smaller strip and
\[
 \|K(t)\|\le Ce^{-\gamma' t}.
\]
\end{theorem}

\begin{proof}
Lift each matrix element to the A2 transfer bundle.  The spectral projector
gives the finite pole part and the Dolgopyat estimate gives the regular
resolvent on the left strip.  A zero of \(\det\widehat C\) would give a
nonzero resolved vector orthogonal to its entire future cyclic span.  Taking
boundary values at the imaginary axis would force zero spectral density and
therefore zero asymptotic variance, contradicting the A2 cohomological
positivity theorem.  Compactness then gives a uniform inverse bound away from
large frequencies.  At large frequency, repeated integration by parts in
\(PU_tP\) uses \(V\subset D(L^r)\) and the A2 polynomial loss, yielding the
displayed estimates.  They are of order strictly below \(|\Im z|^{-2}\) after
choosing \(r>A+4\).  Bromwich inversion may therefore be shifted to
\(\Re z=-\gamma'\); the horizontal sides vanish and give exponential decay.
\end{proof}

\subsection{The genuine Doob tower}

Define
\[
 \widetilde P_t^\Psi F
 =e^{-\lambda_\Psi t}h_\Psi^{-1}P_t^\Psi(h_\Psi F).
\]
By Theorem~\ref{thm:r5-a4-history-feller} this is a Markov semigroup on
\(\mathrm{BUC}_\beta\).  For \(\theta\ne0\), set
\[
 \mathcal E_t^{\Psi,\theta}F
 =\theta^{-1}\log\widetilde P_t^\Psi(e^{\theta F}).
\]

\begin{theorem}[Exact nonlinear history tower]
\label{thm:r5-a4-doob}
The family \(\mathcal E_t^{\Psi,\theta}\) is monotone, cash additive and
satisfies
\[
 \mathcal E_{t+s}^{\Psi,\theta}
 =\mathcal E_t^{\Psi,\theta}\mathcal E_s^{\Psi,\theta}.
\]
Moreover
\[
 D\mathcal E_t^{\Psi,\theta}(0)F
 =\widetilde P_t^\Psi F.
\]
Neither statement uses the generally different ratio
\(P_t^\Psi F/P_t^\Psi\mathbf1\).
\end{theorem}

\begin{proof}
The eigenfunction equation makes \(\widetilde P_t^\Psi\mathbf1=1\), and the
semigroup property follows by cancellation of the middle factors
\(h_\Psi h_\Psi^{-1}\).  Applying the logarithmic exponential transform gives
the nonlinear tower and cash additivity.  Differentiation at zero gives the
last identity.
\end{proof}

\subsection{Enhanced invariance principle with area anomaly}

Let \(G:E\to\mathbb R^d\) be centered and dynamically H\"older.  Define
\[
 D=\int_{-\infty}^{\infty}
 \operatorname{Cov}_\nu(G,G\circ\Theta_t)\,dt,
\]
\[
 \Gamma=\frac12\int_0^\infty
 \left[
  \operatorname{Cov}_\nu(G,G\circ\Theta_t)
 -\operatorname{Cov}_\nu(G\circ\Theta_t,G)
 \right]dt.
\]
Both integrals converge by the A2 spectral gap.

\begin{theorem}[Prepared rough diffusion tangent]
\label{thm:r5-a4-rough}
For
\[
 X_t^\varepsilon=\sqrt\varepsilon
 \int_0^{t/\varepsilon}G(\Theta_s\omega)\,ds
\]
the canonical lift \((X^\varepsilon,\mathbb X^\varepsilon)\) converges in
\(p\)-variation, \(p>2\), to
\[
 B_t,\qquad
 \mathbb B_{s,t}=\int_s^t(B_r-B_s)\otimes\circ dB_r+(t-s)\Gamma,
\]
where \(B\) has covariance \(D\).  Hence the slow equation driven by vector
fields \(V_1,\ldots,V_d\) has limiting drift
\[
 V_0+\sum_{i,j}\Gamma^{ij}[V_i,V_j]
\]
in addition to the Stratonovich diffusion term.
\end{theorem}

\begin{proof}
Solve the Poisson equation in the A2 spectral space and write
\(G=M+\chi-\chi\circ\Theta_1\), where \(M\) is a reverse martingale
difference.  The martingale rough invariance principle gives the Brownian
second level.  Expanding the iterated integral shows that the coboundary terms
vanish except for the deterministic antisymmetric correlation sum, which
converges to \(t\Gamma\).  Exponential correlation decay gives the required
uniform moment bounds and tightness in \(p\)-variation.  The universal limit
theorem for rough differential equations then produces the Lie-bracket drift.
No reversibility assumption is used; when reversibility happens to hold,
\(\Gamma=0\).
\end{proof}

\subsection{The coupled slow--memory limit}

For an integrable memory kernel, write
\[
 K_\varepsilon(t)=\varepsilon^{-1}K(t/\varepsilon),
 \qquad \bar K=\int_0^\infty K(s)\,ds.
\]
The complete prelimit system is
\[
 \begin{aligned}
 dY_t^\varepsilon={}&V_0(Y_t^\varepsilon)dt
 +V(Y_t^\varepsilon)dX_t^\varepsilon
 +B(Y_t^\varepsilon)Z_t^\varepsilon dt,\\
 Z_t^\varepsilon={}&Z_0^\varepsilon(t)
 +\int_0^tK_\varepsilon(t-s)C(Y_s^\varepsilon)\,ds,
 \end{aligned}
\]
where the initial unresolved tail satisfies
\(\sup_{t\le T}|Z_0^\varepsilon(t)|\to0\).

\begin{theorem}[Short-memory and conditional semigroup convergence]
\label{thm:r5-a4-short-memory}
The pair \((Y^\varepsilon,Z^\varepsilon)\) converges to the area-corrected
rough diffusion with
\[
 Z_t=\bar K C(Y_t).
\]
For exposed phases the convergence holds for conditional transition kernels,
uniformly on compact sets of regular histories.  Consequently the exact Doob
history semigroups converge locally uniformly to the Markov diffusion
semigroup of the corrected slow equation.
\end{theorem}

\begin{proof}
Theorem~\ref{thm:r5-a4-memory-decay} makes
\(K\) exponentially integrable.  Tightness of \(Y^\varepsilon\) and uniform
continuity of \(C\) give
\[
 \sup_{t\le T}\left|
  \int_0^tK_\varepsilon(t-s)C(Y_s^\varepsilon)ds
  -\bar K C(Y_t^\varepsilon)
 \right|\to0
\]
in probability; split the integral at \(t-\varepsilon^{1/2}\) and use the
exponential tail on the old part.  Combine this with
Theorem~\ref{thm:r5-a4-rough} and rough-path stability.

For conditional kernels, couple two histories with the same present resolved
coordinate until their common magnet return.  The history Feller estimate and
the exponential memory tail bound the difference of future cylinder
expectations by
\(Ce^{-\gamma m}+C\varepsilon^\alpha\) when their last \(m\) units agree.
First let \(\varepsilon\to0\), then \(m\to\infty\).  This proves convergence
of optional projections and hence of the Doob kernels, not merely convergence
under one stationary initial law.
\end{proof}

\begin{theorem}[Round-five A4 closure]
\label{thm:r5-a4-main}
Prepared paths possess a precise measured Ray realization; exposed Sinai
phases have an explicit norm-Feller history kernel and positive Doob
eigenfunction; finite resolved coordinates have a domain-safe exponentially
decaying Volterra memory; and their deterministic homogenization limit
includes the generally nonzero rough-area anomaly and the complete
short-memory coupling.
\end{theorem}

\begin{proof}
Combine Theorems~\ref{thm:r5-a4-ray},
\ref{thm:r5-a4-history-feller}, \ref{thm:r5-a4-volterra},
\ref{thm:r5-a4-memory-decay}, \ref{thm:r5-a4-doob},
\ref{thm:r5-a4-rough}, and \ref{thm:r5-a4-short-memory}.
\end{proof}
